\label{sec:ifc-tu}

Any C++ translation unit can be compiled into an IFC structure.  A translation unit so represented can be referenced by abstract reference
of type \type{UnitIndex}.  
%
\begin{figure}[htbp]
	\centering
	\absref{3}{UnitSort}
	\caption{\code{UnitIndex}: Abstract reference of a declaration}
	\label{fig:ifc-unit-index}
\end{figure}

\begin{SortEnum}{UnitSort}
	\enumerator{Source}
	\enumerator{Primary}
	\enumerator{Partition}
	\enumerator{Header}
	\enumerator{ExportedTU}
	\enumerator{Archive}
\end{SortEnum}

with meaning as explained in the sections below.  The \field{index} value has a tag-dependent interpretation as defined below.

\section{Translation unit structures}
\label{sec:ifc:tu-structures}

\subsection{\valueTag{UnitSort::Source}}
\label{sec:ifc:UnitSort:Source}
A \type{UnitIndex} value with this tag designates a translation unit defined by a general C++ source file.

The \field{index} value is undefined and has no meaning.


\subsection{\valueTag{UnitSort::Primary}}
\label{sec:ifc:UnitSort:Primary}

A \type{UnitIndex} value with this tag designates a primary module interface.

The \field{index} is to be interpreted as a \type{TextOffset} value (\secref{sec:ifc-textoffset-data-type}), 
which is an offset into the string table (\secref{sec:ifc-string-table}).  
The string at that offset is the name of the module (\secref{sec:ifc-modules}) for which this translation unit is a primary module interface.


\subsection{\valueTag{UnitSort::Partition}}
\label{sec:ifc:UnitSort:Partition}

A \type{UnitIndex} value with this tag designates a partition module unit.

The \field{index} is to be interpreted as a \type{TextOffset} value that designates the name of the module partition.  
The string of that name is \texttt{M:P}, obtained as the concatenation of the name \texttt{M} of the parent module, 
the colon character (\texttt{:}), and the relative name \texttt{P} of the partition.

\subsection{\valueTag{UnitSort::Header}}
\label{sec:ifc:UnitSort:Header}

A \type{UnitIndex} value with this tag designates a C++ header unit.  
The \field{index} value is to be interpreted as a \type{TextOffset} value (\secref{sec:ifc-textoffset-data-type})
into the string table (\secref{sec:ifc:string-table}); the string at that offset is the name of the
header unit, \ie{} the \grammar{header-name} as specified when IFC for the header unit was built.


\subsection{\valueTag{UnitSort::ExportedTU}}
\label{sec:ifc:UnitSort:ExportedTU}

A \type{UnitIndex} value with this tag designates a translation unit compiled by the MSVC compiler with the compiler flags
\texttt{/module:export} and \texttt{/module:name}.  Such a translation is processed as if every toplevel declaration was
prefixed with the keyword \texttt{export}.  This is an MSVC extension.

The \field{index} is to be interpreted as a \type{TextOffset} value that designates the name specified via the compiler
switch \texttt{/module:name}.

\note{An IFC unit of this sort is deprecated and scheduled for removal from MSVC.}

\subsection{\valueTag{UnitSort::Archive}}
\label{sec:ifc:UnitSort:Archive}

A \type{UnitIndex} value with this tag designates an IFC file that is an \emph{archive}: a conglomerate that packages
a set of other, non-archive, IFC files.  An archive is useful as a ``packaging'' container for smaller IFC files, not unlike
object file archives (\eg{} \code{.a} files).

The \field{index} is to be interpreted  as a \type{TextOffset} value (\secref{sec:ifc-textoffset-data-type}) into the
archive's string table (\secref{sec:ifc-string-table}).  When non-null, the string at that offset is a name for the archive
as a whole; the null offset (value $0$) means the archive is unnamed.

Each packaged IFC file is a \emph{member} of the archive.  A member is identified by its
\emph{sort} (its own \type{UnitSort}) together with is \emph{canonical name}, the name
recorded in the member's own \field{unit} descriptor (\secref{sec:ifc-tu}) -- \eg{} the module name of a primary
module interface (\valueTag{UnitSort::Primary}) or module partition (\valueTag{UnitSort::Partition}), or the
header-name of a header unit (\valueTag{UnitSort::Header}).  Members of different sorts may therefore share
a canonical name, \eg{} a named module and a header unit both named \code{M}.  A member's content
is the complete IFC image (file signature, header, partitions, string table, and table of contents), stored verbatim.

A well-formed archive IFC satisfies the following constraints:
\begin{itemize}
\item Every member has a canonical name; an IFC whose \field{unit} descriptor names no unit 
	(\eg{} \valueTag{UnitSort::Source} with no name) cannot be a member.
\item  No two members share the same sort and the same canonical name.
\item The \type{ArchiveMember} records (\secref{sec:ifc-archive-member}) are sorted by canonical name,
	with the member's sort distinguishing members that share a name, so the table of contents can be 
	binary searched by name.
\item Members are non-archive IFCs: an archive cannot contain another archive (\ie{} no \type{ArchiveMember} has a member whose
	\field{tag} is also \valueTag{UnitSort::Archive}).
\end{itemize}

In all other respects an archive IFC is an ordinary IFC file: it begins with the file signature (\secref{sec:ifc-file-signature});
its header (\secref{sec:ifc-file-header}) is covered by the content hash (\secref{sec:ifc-content-sha256-hash}) as usual; and the
header, string table, and table of contents are stored in little-endian (\secref{sec:ifc-endianness}).  
Each begins at a 4-byte boundary.
